\ExerciseSection

\begin{enumerate}
\def\labelenumi{\roman{enumi})}
\tightlist
\item
  \begin{enumerate}
  \def\labelenumii{\alph{enumii})}
  \tightlist
  \item
    构造一个由四个点组成的拓扑空间，它满足公理 \((\mathbf{O}_{\mathbf{V}})\) 但不满足公理 \((\mathbf{O}_{\mathbf{III}})\) 。
  \end{enumerate}
\end{enumerate}

\begin{enumerate}
\def\labelenumi{\alph{enumi})}
\setcounter{enumi}{1}
\item
  满足公理 \((\mathrm{O}_{\mathrm{V}})\) 的拓扑空间 X 也满足公理 \((\mathrm{O}_{\mathrm{III}})\) ，当且仅当它具有下列性质：X 的每个闭子集都是它的诸邻域的交。这时 X 是可一致化的，并且与 X 相伴的完全正则空间（§ 1，习题 4）是正规的。
\item
  若拓扑空间满足公理 (C) 与 \((\mathrm{O}_{\mathrm{III}})\) ，则它满足 \((\mathrm{O}_{\mathrm{V}})\) ，且相伴的完全正则空间是紧的。
\end{enumerate}

\begin{enumerate}
\def\labelenumi{\arabic{enumi})}
\setcounter{enumi}{1}
\tightlist
\item
  设 X 是满足公理 \((\mathbf{O}_{\mathbf{v}})\) 的拓扑空间。
\end{enumerate}

\begin{enumerate}
\def\labelenumi{\alph{enumi})}
\item
  证明 X 的两点 x、y 之间的关系 \(R: \{x\} \cap \{y\} \neq \emptyset\) 等价于关系 \(R' : "f(x) = f(y)\) ——对 X 上每个实值连续函数 f——''；由此证明 R 是 X 上的一个等价关系。
\item
  证明商空间 \(\mathbf{Y} = \mathbf{X} / \mathbf{R}\) 是正规的，且若 \(\varphi: \mathbf{X} \to \mathbf{Y}\) 是典范映射，则每个连续实值函数 \(f\) （在 \(\mathbf{X}\) 上）都可写成 \(f = g \circ \varphi\) 的形式，其中 \(g\) 是 \(\mathbf{Y}\) 上的一个连续实值函数。
\end{enumerate}

\begin{enumerate}
\def\labelenumi{\arabic{enumi})}
\setcounter{enumi}{2}
\tightlist
\item
  若 X 是一个拓扑空间，则下列条件等价：
\end{enumerate}

\begin{enumerate}
\def\labelenumi{\alph{enumi})}
\item
  X 的每个子空间满足公理 \((\mathrm{O}_{\mathrm{V}}^{\prime})\) 。
\item
  X 的每个开子空间满足公理 \((\mathbf{O}_{\mathbf{V}}^{\prime})\) 。
\item
  对 X 的任意两个满足 \(A \cap \overline{B} = B \cap \overline{A} = \varnothing\) 的子集 A 与 B，存在两个不相交的开集 U、V，使得 \(A \subset U\) 且 \(B \subset V\) 。
\end{enumerate}

拓扑空间 \(\mathbf{X}^{-}\) 称为完全正规的，如果它是豪斯多夫的且满足这些条件。每个可度量化空间都是完全正规的。紧空间未必是完全正规的。

\begin{enumerate}
\def\labelenumi{\arabic{enumi})}
\setcounter{enumi}{3}
\tightlist
\item
  每个全序集 X，赋以两个拓扑 \(\mathcal{C}_{+}(\mathbf{X})\) 、\(\mathcal{C}_{-}(\mathbf{X})\) 之一（第 I 章，§ 2，习题 5），都是完全正规的。（若 A 与 B 是 X 的两个满足 \(\mathbf{A} \cap \overline{\mathbf{B}} = \mathbf{B} \cap \overline{\mathbf{A}} = \emptyset\) 的子集，定义邻域 \(V_{x}\) （\(x\) 的，对每个 \(x \in \mathbf{A}\) ），以及邻域 \(W_{y}\) （\(y\) 的，对每个 \(y \in \mathbf{B}\) ），使得 \(V_{x} \cap W_{y} = \emptyset\) 对所有 \(x \in \mathbf{A}\) 及所有 \(y \in \mathbf{B}\) 成立。）
\end{enumerate}

¶ 5) 每个全序集 X，赋以拓扑 \(\mathcal{G}_{0}(\mathbf{X})\) （第 I 章，§ 2，习题 5），都是完全正规的。{[}先考虑 X 是紧的情形；借助第 IV 章，§ 2，习题 6，证明 X 中每个开集都是两两不相交的开区间之并；利用这一结果证明命题，办法是（用习题 3 的记号）考虑闭集 \(\overline{\mathbf{A}}\cap\overline{\mathbf{B}}\) 的余集，然后考虑 \(\overline{\mathbf{B}}\) 的余集；为过渡到 X 任意的普遍情形，利用第 IV 章，§ 4，习题 7。{]}

\begin{enumerate}
\def\labelenumi{\arabic{enumi})}
\setcounter{enumi}{5}
\item
  证明：在正规空间 X 中，X 的每个是闭集的可数并的子空间 Y 都是正规的。{[}设 A、B 是 Y 的两个不相交的闭子集，并设 Y 是闭子集组成的一个递增序列 \((\mathbf{Y}_n)\) （在 \(\mathbf{X}\) 中）的并。用归纳法定义开子集的两个序列 \((\mathbf{U}_n), (\mathbf{V}_n)\) （\(\mathbf{X}\) 的），使得：(i) \(\overline{\mathbf{U}_n \cap \mathbf{Y}_n}\) 与 \(\overline{\mathbf{V}_n \cap \mathbf{Y}_n}\) 不相交；(ii) \(\mathbf{U}_n \cap \mathbf{Y}_n\) 包含 \(\mathbf{A} \cap \mathbf{Y}_n\) 以及诸集 \(\overline{\mathbf{U}_i \cap \mathbf{Y}_i}\) （对 \(1 \leqslant i \leqslant n\) ）；且 \(\mathbf{V}_n \cap \mathbf{Y}_n\) 包含 \(\mathbf{B} \cap \mathbf{Y}_n\) 以及诸 \(\overline{\mathbf{V}_i \cap \mathbf{Y}_i}, 1 \leqslant i \leqslant n\) 。{]}
\item
  拓扑空间 X 称为完备正规的，如果它是正规的，且 X 的每个闭子集都是可数个开集的交（或者等价地，X 的每个开子集都是可数个闭集的并）。
\end{enumerate}

\begin{enumerate}
\def\labelenumi{\alph{enumi})}
\item
  豪斯多夫空间 X 是完备正规的，当且仅当对 X 的任一闭子集 A，存在 X 上的一个实值连续函数 f，使得 \(\overline{f}^{1}(0)=A\) {[}利用 § 1，习题 16 a) 的方法{]}。b) 证明每个完备正规空间 X 都是完全正规的，且 X 的每个子空间都是完备正规的{[}利用习题 3 b) 与 6{]}。
\item
  完备正规空间上的每个下半连续实值函数都是连续函数组成的一个序列的上包络（利用 § 2，第 7 号，命题 11 的方法）。
\item
  在不可数离散空间上加一个无穷远点所得的紧空间 X 是完全正规但非完备正规的，且 X 的每个子空间都是仿紧的。
\end{enumerate}

§ 8) 证明 § 2，习题 13 a) 中定义的非可度量化紧空间 X 是完备正规的，而 § 2，习题 13 d) 中定义的非可度量化紧空间 Y 是完全正规的但不是完备正规的。

¶ 9) a) 设 X 是一个仿紧空间，其中每个点都有可数的基本邻域系，Y 是一个可数紧的正规空间（§ 2，习题 14）。证明 X × Y 是正规的。{[}设 A、B 是 X × Y 的两个不相交的闭子集；对每个 x ∈ X，证明存在 x 在 X 中的一个邻域 U\_x 及 Y 中两个开集 V\_x、W\_x，使得 \(\overline{V}_x \cap \overline{W}_x = \emptyset\) ，且对每个 z ∈ U\_x 有 A(z) ⊂ V\_x 及 B(z) ⊂ W\_x；为证这一点，用反证法论证。设 (T\_λ)\_\{λ ∈ L\} 是 X 的一个局部有限开覆盖，它加细覆盖 (U\_x)\_\{x ∈ X\}，并设 (f\_λ) 是从属于 (T\_λ) 的一个单位分解；对每个 λ ∈ L，取 x\_λ 使 T\_λ ⊂ U\_\{x\_λ\}，并令 g\_λ 为 Y 到 {[}0, 1{]} 中的一个连续映射，它在 \(\overline{V}_{x_λ}\) 上等于 0，在 \(\overline{W}_{x_λ}\) 上等于 1；考虑 X × Y 上的函数 \(\sum_{\lambda} f_\lambda(x) g_\lambda(y)\) 。{]}

\begin{enumerate}
\def\labelenumi{\alph{enumi})}
\setcounter{enumi}{1}
\tightlist
\item
  设 Y = {[}a, b{[} 是第 I 章，§ 9，习题 12 中定义的局部紧空间，X = {[}a, b{]} 是在 Y 上加一个无穷远点所得的紧空间。Y 是正规的（习题 4）且可数紧（§ 2，习题 15），但乘积 X×Y 不是正规的（利用第 II 章，§ 4，习题 4）。
\end{enumerate}

¶ 10) 设 L 是一个不可数集，X 表示完全正则空间 \(\mathbf{N}^{\mathrm{L}}\) （N 赋以离散拓扑）。设 A（相应地 B）是 X 中所有满足如下条件的 \(x = (x(\lambda))_{\lambda \in \mathrm{L}}\) 组成的集合：对每个整数 \(k \neq 0\) （相应地 \(k \neq 1\) ），由 \(\lambda \in \mathbf{L}\) 中满足 \(x(\lambda) = k\) 者组成的集至多有一个元素。a) 证明 A 与 B 是 X 的两个不相交的闭子集。

\begin{enumerate}
\def\labelenumi{\alph{enumi})}
\setcounter{enumi}{1}
\item
  设 U、V 是 X 的两个开子集，满足 \(A \subset U\) 与 \(B \subset V\) 。令 \(\Phi\) 表示 X 的所有初等集（第 I 章，§ 4，第 1 号）中其异于 N 的投影均由单独一个元素组成者所成的集合；对每个 \(W \in \Phi\) ，令 \(\mathrm{H}(\mathrm{W})\) 为所有满足下列条件的 \(\lambda \in L\) 组成的集合：\(\mathrm{pr}_{\lambda}(\mathrm{W})\) 由单独一个元素组成。证明存在 A 的点组成的一个序列 \((x_{n})\) 、一个集序列 \((\mathbf{U}_{n})\) （集取自 \(\Phi\) ）、L 的不同元素组成的一个序列 \((\lambda_{n})\) 以及整数的一个严格递增序列 \((m(n))_{n \in \mathbb{N}}\) ，它们具有下列性质：(i) \(U_{n} \subset U\) 是 \(x_{n}\) 的一个邻域；(ii) \(\mathrm{H}(\mathrm{U}_{n})\) 是诸 \(\lambda_{k}\) （\(k \leqslant m(n)\) ）组成的集合；(iii) \(x_{0}(\lambda) = 0\) 对每个 \(\lambda \in L\) 成立，且对每个 n \textgreater{} 0，\(x_{n}(\lambda_{k}) = k\) （若 \(k \leqslant m(n - 1)\) ），而 \(x_{n}(\lambda) = 0\) 对所有其他 \(\lambda \in L\) 成立。
\item
  设 \(y \in \mathbf{B}\) 是这样的点：\(y(\lambda_k) = k\) 对每个整数 \(k\) 成立，且 \(y(\lambda) = \mathbf{1}\) 对所有 \(\lambda \in \mathbf{L}\) 中异于诸 \(\lambda_k\) 者成立。设 \(\mathbf{V}_0 \subset \mathbf{V}\) 是 \(\Phi\) 的一个包含 \(y\) 的集合，并设 \(n\) 是使 \(\lambda_k \in \mathbf{L} - \mathbf{H}(\mathbf{V}_0)\) （对所有 \(k > m(n)\) ）的整数。证明 \(\mathbf{U}_{n+1} \cap \mathbf{V}_0 \neq \emptyset\) ，并推出 \(\mathbf{X}\) 不是正规的。
\end{enumerate}

\begin{enumerate}
\def\labelenumi{\arabic{enumi})}
\setcounter{enumi}{10}
\tightlist
\item
  由习题 10 推出：
\end{enumerate}

\begin{enumerate}
\def\labelenumi{\alph{enumi})}
\item
  若非空豪斯多夫空间的一个乘积 \(\prod_{i\in I}X_i\) 是正规的，则除去一个可数指标集外，\(X_i\) 都是可数紧的（§ 2，习题 14）。
\item
  非空可度量化空间的乘积 \(\prod_{i\in I}X_i\) 是正规的，当且仅当除去一个可数指标集外 \(X_{i}\) 都是紧的；此时乘积空间 \(\prod_{i\in I}X_{i}\) 是仿紧的。
\end{enumerate}

\begin{enumerate}
\def\labelenumi{\arabic{enumi})}
\setcounter{enumi}{11}
\tightlist
\item
  设 X、Y 是两个豪斯多夫空间。
\end{enumerate}

\begin{enumerate}
\def\labelenumi{\alph{enumi})}
\item
  假设 X 中有一个闭集 A，它不是可数个开集的交，且 Y 有一个可数无限子集 B，它不是闭的。取 \(b \in \overline{B} - B\) ，并考虑子集 \(C = A \times B\) 与 \(D = (X - A) \times \{b\}\) （在 \(X \times Y\) 中）。证明 \(C \cap \overline{D} = \overline{C} \cap D = \varnothing\) ，但不存在 \(X \times Y\) 的开子集的对 U、V，使得 \(C \subset U\) 、\(D \subset V\) 且 \(U \cap V = \varnothing\) 。
\item
  若 \(X \times Y\) 是完全正规的，则或者 (i) 空间 X、Y 之一是完备正规的，或者 (ii) 空间 X、Y 之一具有每个可数无限子集都闭的性质（参看 § 5，习题 16）。
\item
  设 \(X = [a, b]\) 是一个不可数的良序集，使得 \([a, x]\) 对每个 x \textless{} b 都是可数的。给 X 赋以如下拓扑：\(\{x\}\) 对每个 x \textless{} b 是开的，且诸区间 \([x, b]\) （x \textless{} b）构成 b 的一个基本邻域系。证明 \(X \times X\) 是完全正规的，但 X 不是完备正规的。
\item
  证明：若 \(\mathbf{X}\) 是紧空间，且 \(\mathbf{X} \times \mathbf{X} \times \mathbf{X}\) 完全正规，则 \(\mathbf{X}\) 可度量化（利用 § 2，第 4 号，定理 1）。
\end{enumerate}

¶ 13) 设 \((\mathbf{X}_{t})_{t\in I}\) 是豪斯多夫拓扑空间组成的一个不可数族，其中每个空间都不止含一个点。对每个 \(t\in I\) ，设 \(a_{t}\) 与 \(b_{t}\) 是 \(\mathbf{X}_{t}\) 的两个不同的点。设 \(\mathbf{X}\) 是积空间 \(\prod_{t\in I}\mathbf{X}_{t}\) ，并设 Y 是 X 中由这样的点 \((x_{t})_{t\in I}\) 组成的子空间：除去一个可数的指标集外，均有 \(x_{t}=a_{t}\) ；又设 b 是 X 的点 \(b_{t}\) 。

在积空间 \(\mathbf{X} \times \mathbf{Y}\) 中，设 \(\mathbf{A}\) 是 \(\mathbf{Y} \times \mathbf{Y}\) 的对角线，并设 \(\mathbf{B}\) 是集 \(\{b\} \times \mathbf{Y}\) 。证明 \(\mathbf{A}\) 与 \(\mathbf{B}\) 是闭集，并且不存在开集对 \(\mathbf{U}, \mathbf{V}\) （在 \(\mathbf{X} \times \mathbf{Y}\) 中），满足 \(\mathbf{A} \subset \mathbf{U}, \mathbf{B} \subset \mathbf{V}\) 以及 \(\mathbf{U} \cap \mathbf{V} = \emptyset\) 。{[}设 \(\mathbf{U}, \mathbf{V}\) 是两个开集，满足 \(\mathbf{A} \subset \mathbf{U}\) 与 \(\mathbf{B} \subset \mathbf{V}\) ；证明存在一个递增序列 \((\mathbf{H}_n)\) （由 \(\mathbf{I}\) 的可数子集组成），使得：若 \(x_n\) 表示 \(\mathbf{Y}\) 中这样的点——其指标 \(i \in \mathbf{H}_n\) 的每个坐标等于 \(b_i\)，而其指标 \(i \notin \mathbf{H}_n\) 的每个坐标等于 \(a_i\) ——则点 \((x_{n+1}, x_n)\) 属于 \(\mathbf{V}\) （对每个整数 \(n\) ）；证明序列 \((x_{n+1}, x_n)\) 趋于 \(\mathbf{A}\) 的一个点，并由此推出 \(\mathbf{U} \cap \mathbf{V} \neq \emptyset\) 。{]}

由此构造一个非正规的连通拓扑环的例子。

并由此推出：由豪斯多夫拓扑空间组成的一个不可数族（其中每个空间至少有两个不同的点）的积不可能是完全正规的（在诸 \(X_{t}\) 都是两个点的离散空间时利用上述结果）。

\begin{enumerate}
\def\labelenumi{\arabic{enumi})}
\setcounter{enumi}{13}
\item
  设 X 是一个非正规的豪斯多夫空间，并设 A, B 是 X 的两个不相交的闭子集，使得不存在 X 中满足关系 \(A \subset U\) 与 \(B \subset V\) 的不相交开集对 U, V。设 R 是 X 上的等价关系，其等价类为：集 A、集 B、以及诸集 \(\{x\}\) （其中 \(x \in \complement(A \cup B)\) ）。证明 R 的图象在 \(X \times X\) 中是闭的，且关系 R 是闭的，但商空间 X/R 不是豪斯多夫空间（参看第 I 章，§ 8，第 3 号，命题 8）。
\item
  \begin{enumerate}
  \def\labelenumii{\alph{enumii})}
  \tightlist
  \item
    设 X 是正规空间，R 是 X 上的一个闭等价关系。证明商空间 X/R 是正规的（利用第 I 章，§ 5，第 4 号，命题 10）。
  \end{enumerate}
\end{enumerate}

\begin{enumerate}
\def\labelenumi{\alph{enumi})}
\setcounter{enumi}{1}
\item
  设 X 是一个局部紧的 σ 紧空间，R 是 X 上的等价关系，其图象在 X×X 中是闭的。证明 X/R 是正规的（参看第 I 章，§ 10，习题 19）。
\item
  设 X 是 R 中由形如 1/n（其中 n 是异于 0 与 ±1 的整数）的点组成的集的余集。在可度量化空间 X 上考虑这样的等价关系 R：每个非整数的 \(x \in X\) 的等价类由点 x 与 1/x 组成，而每个整数的等价类由该整数单独组成。证明 R 是开的，且 R 的图象在 \(X \times X\) 中是闭的，但 X/R 不是正则的。
\end{enumerate}

\begin{enumerate}
\def\labelenumi{\arabic{enumi})}
\setcounter{enumi}{15}
\tightlist
\item
  对每个覆盖 \(\Re\) （拓扑空间 \(X\) 的），用 \(V_{\Re}\) 表示诸集 \(U \times U\) （\(U \in \Re\)）的并。覆盖 \(\Re\) （\(X\) 上的）称为均匀的，如果有一个邻域 \(W\) （对角线 \(\Delta\) 的，\(X \times X\) 中的），使得覆盖 \((W(x))_{x \in X}\) 加细 \(\Re\)。称 \(\Re\) 为可除的，如果有一个邻域 \(W\) （\(\Delta\) 的，在 \(X \times X\) 中），使得 \(\dot{W} \subset V_{\Re}\) 。 a) 证明 \(X\) 的每个均匀开覆盖都是可除的{[}参看习题 19 b){]}。
\end{enumerate}

\begin{enumerate}
\def\labelenumi{\alph{enumi})}
\setcounter{enumi}{1}
\tightlist
\item
  设 \(\Re\) 是 \(X\) 的一个开覆盖，使得存在一个局部有限闭覆盖 \(\mathfrak{S}\) （\(X\) 的）加细 \(\Re\)。证明 \(\Re\) 是均匀的。（对每个集 \(A \in \mathfrak{S}\)，设 \(U_A\) 是 \(\Re\) 中包含 \(A\) 的一个集，并设 \(V_A\) 是 \(X \times X\) 中 \(U_A \times U_A\) 与 \((X - A) \times (X - A)\) 的并；考虑集 \(W = \bigcap_{A \in \mathfrak{S}} V_A\) 。）
\end{enumerate}

¶ 17) a) 若拓扑空间 X 的每个有限开覆盖都可除（习题 16），则 X 满足公理 \((\mathrm{O}_{\mathbf{V}}^{\prime})\) （考虑 X 的一个二元开覆盖，即由 X 的两个子集所组成的覆盖）。反之，若 X 满足 \((\mathrm{O}_{\mathbf{V}}^{\prime})\) ，则 X 的每个有限开覆盖都是均匀的{[}利用第 3 号的定理 3（只要 X 满足 \((\mathrm{O}_{\mathbf{V}}^{\prime})\) ，它便成立），先化归到二元覆盖的情形，然后再处理这一情形{]}。当 R 跑遍 X 的有限开覆盖所成的集时，诸集 \(V_{R}\) 构成 X 上一个一致结构的围域的基本系，该一致结构与 X 的拓扑相容；这个一致结构称为``有限开覆盖的一致结构''。

\begin{enumerate}
\def\labelenumi{\alph{enumi})}
\setcounter{enumi}{1}
\item
  假设 X 是正规的。证明有限开覆盖的一致结构 U 与如下的最粗一致结构 \(U'\) 相同：X 到 {[}0, 1{]} 的所有连续映射关于它都一致连续{[}即 X 上由其斯通–切赫紧化（§ 1，习题 7）的一致结构所诱导的一致结构{]}。{[}为证 U 比 \(U'\) 粗，利用第 3 号的定理 3 与公理 \((\mathrm{O_V})\)，在 \(\mathbf{X}\) 上构造形如 \((x, y) \to |f(x) - f(y)|\) 的伪度量组成的一个有限族，使得由这些伪度量所定义的一致结构的一个围域含于一致结构 \(\mathfrak{U}]\) 的一个围域。
\item
  设 X 是正规空间，并设 \(\beta\mathbf{X}\) 是它的斯通–切赫紧化。对 X 的每个闭子集 Y，设 \(\overline{\mathbf{Y}}\) 是 Y 在 \(\beta\mathbf{X}\) 中的闭包。证明由 Y 的恒等映射扩张而得的、从 \(\beta\mathbf{Y}\) 到 \(\overline{\mathbf{Y}}\) 的连续映射（\(\beta\mathbf{Y}\) 是 Y 的斯通–切赫紧化）是一个同胚{[}或者利用 \(b\) )，或者利用定理 2 与 \(\beta\mathbf{Y}]\) 的万有性质。若 Y 与 Z 是 X 的两个闭子集，证明 \(\overline{\mathbf{Y}} \cap \overline{\mathbf{Z}} = \overline{\mathbf{Y} \cap \mathbf{Z}}\) 在 \(\beta\mathbf{X}\) 中成立（若 \(x_0 \notin \overline{\mathbf{Y}} \cap \overline{\mathbf{Z}}\) 且 \(x_0 \in \overline{\mathbf{Y}}\) ，考虑 \(x_0\) 在 \(\beta\mathbf{X}\) 中的一个闭邻域 V，使得 V \(\cap\) Y 与 Z 是 X 的不相交闭子集）。
\end{enumerate}

¶ 18) 拓扑空间 X 的子集组成的一个族 \((A_{\alpha})\) 称为离散的，如果对每个 \(x \in X\) ，存在 x 的一个邻域，它至多与族中的一个集 \(A_{\alpha}\) 相交。豪斯多夫空间 X 称为集体正规的，如果对 X 的闭子集组成的每个离散族 \((A_{\alpha})\) ，存在 X 的两两不相交的开子集组成的一个族 \((U_{\alpha})\) ，使得 \(A_{\alpha} \subset U_{\alpha}\) 对每个指标 \(\alpha\) 成立。每个集体正规空间都是正规的。

\begin{enumerate}
\def\labelenumi{\alph{enumi})}
\item
  完全正则空间 X 的每个开覆盖都可除（习题 16），当且仅当对角线 \(\Delta\) （在 \(\mathbf{X} \times \mathbf{X}\) 中）的邻域所成之集是 X 的万有一致结构的围域滤子（ \(\S\) 1，习题 5）。证明：若此条件满足，则 X 是集体正规的{[}若 \((\mathbf{A}_{\alpha})\) 是 X 的闭子集组成的一个离散族，考虑开覆盖 \((\mathbf{V}_{\alpha})\) ，其中 \(\mathbf{V}_{\alpha} = \mathbf{X} - \bigcup_{\beta \neq \alpha} \mathbf{A}_{\beta}\) 对每个 \(\alpha\) 成立{]}。
\item
  设 Y = {[}a, b{]} 是第 I 章，§ 9，习题 12 中定义的局部紧空间；设 Y₀ 是赋以离散拓扑的集 Y，并设 Z 是集 Y₀ ∪ \{b\}，其中 Y₀ 的点是开集，而诸集 {]}x, b{[}（x 跑遍 Y₀）构成 b 的一个基本邻域系。证明空间 X = Y × Z 是集体正规的（回忆 Y 的子集组成的任何无限族都不可能是离散的，并利用习题 4）。设 R 是 X 的开覆盖，由 Y × Y₀ 以及诸积 {[}a, x{]} × {]}x, b{]} （其中 x \textless{} b）组成；证明 R 不是可除的，并且对角线 Δ （在 X × X 中的）的邻域所成之集因而不是 X 上一个一致结构的围域滤子{[}利用第 I 章，§ 9，习题 12 a){]}。
\end{enumerate}

¶ 19) a) 正则空间 X 是仿紧的，当且仅当 X 的每个开覆盖都是均匀的（习题 16）。{[}为证必要性，利用引理 5 与习题 16 b)；为证充分性，利用 § 4 的习题 16 a)、§ 1 第 4 号的命题 2 以及 § 4 第 5 号的定理 4{]}。

\begin{enumerate}
\def\labelenumi{\alph{enumi})}
\setcounter{enumi}{1}
\item
  给出集体正规但非仿紧的空间上的一个可除而非均匀的覆盖的例子{[}利用 a) 与第 II 章，§ 4，习题 4{]}。
\item
  证明仿紧空间 X 关于其万有一致结构（§ 1，习题 5）是完备的。{[}若关于这个一致结构的一个柯西滤子 F 没有触点，则每个点 \(x \in X\) 都有一个开邻域 \(V_{x}\) ，它至少不与 F 的一个集相交；利用 a) 与习题 18 a)。{]}
\end{enumerate}

\begin{enumerate}
\def\labelenumi{\arabic{enumi})}
\setcounter{enumi}{19}
\tightlist
\item
  \begin{enumerate}
  \def\labelenumii{\alph{enumii})}
  \tightlist
  \item
    正则空间 X 是仿紧的，如果对 X 的每个开覆盖 R，存在一个序列 \((\mathfrak{S}_{n})\) ，使得 \(\mathfrak{S} = \bigcup_{n} \mathfrak{S}_{n}\) 是 \(\mathbf{X}\) 的一个加细 \(R\) 的开覆盖，且每个 \(\mathfrak{S}_n\) 都是局部有限族（参看引理 4、5、6）。
  \end{enumerate}
\end{enumerate}

\begin{enumerate}
\def\labelenumi{\alph{enumi})}
\setcounter{enumi}{1}
\item
  由 a) 推出：在仿紧空间中，闭集组成的每个可数并都是一个仿紧子空间（参看 § 5，习题 15）。
\item
  在正则空间 X 中，设 F 是闭集组成的一个局部有限族，其中每个闭集都是 X 的仿紧子空间。证明 F 中诸集的并是 X 的一个仿紧子空间（利用第 5 号的引理 6；参看 § 5，习题 15）。
\item
  证明：若 X 是仿紧空间，而 Y 是紧集组成的一个可数并的正则空间，则 \(X \times Y\) 是仿紧的{[}把 Y 嵌入它的斯通–切赫紧化并利用 b)；参看 § 5，习题 16{]}。
\end{enumerate}

¶ 21) a) 设 X 是仿紧空间，R 是 X 上的一个闭等价关系，使得关于 R 的每个等价类都是紧的。证明 X/R 是仿紧的{[}利用习题 15 a)、第 3 号的定理 3 以及第 5 号的引理 6；参看习题 15 c){]}。

\begin{enumerate}
\def\labelenumi{\alph{enumi})}
\setcounter{enumi}{1}
\item
  设 X 是豪斯多夫空间，R 是 X 上的一个闭等价关系，使得关于 R 的每个等价类都是紧的。证明：若 X/R 仿紧，则 X 也仿紧（利用第 I 章，§ 9，第 10 号，命题 17 的论证）。
\item
  设 Y 是局部紧但非仿紧的空间（参看第 I 章，§ 9，习题 12），并设 \(Y_{0}\) 是 Y 的一点紧化。设 X 是 Y 与 \(Y_{0}\) 的拓扑和，并设 R 是 X 上的等价关系，它把 Y 的每个点与其在 \(Y_{0}\) 中的典范象等同起来。关系 R 是开的，关于 R 的每个等价类至多含两个点，且 X/R 同胚于 \(Y_{0}\) ，因而是紧的。
\end{enumerate}

\begin{enumerate}
\def\labelenumi{\arabic{enumi})}
\setcounter{enumi}{21}
\item
  正则空间 \(\mathbf{X}\) 可度量化，当且仅当存在一个序列 \((\mathfrak{B}_n)\) （由 \(\mathbf{X}\) 的开子集的局部有限族组成），使得 \(\mathfrak{B} = \bigcup_{n} \mathfrak{B}_n\) 是 \(\mathbf{X}\) 的拓扑的一个基。{[}为证条件是必要的，利用第 5 号的引理 3。为证它是充分的，先用引理 4、5、6 证明 X 是仿紧的。对每对整数 m, n 以及每个 \(U \in \mathfrak{B}_{m}\) ，设 \(U'\) 是诸集 \(V \in \mathfrak{B}_{n}\) （其中 \(\overline{V} \subset U\) ）的并。证明 \(\overline{U}' \subset U\) 。设 \(f_{U}: X \to [0, 1]\) 是一个连续映射，它在 X---U 上等于 0，在 \(U'\) 上等于 1；并设 \(d_{mn}(x, y) = \sum_{U \in \mathfrak{B}_{m}} |f_{U}(x) - f_{U}(y)|\) ；证明诸伪度量 \(d_{mn}\) 定义 X 的拓扑。{]}（``长田–斯米尔诺夫定理''）。特别地，具有可数基的每个正则空间都是可度量化的。
\item
  \begin{enumerate}
  \def\labelenumii{\alph{enumii})}
  \tightlist
  \item
    每个正则的林德勒夫空间（第 I 章，§ 9，习题 15）都是仿紧的{[}利用习题 20 a){]}。
  \end{enumerate}
\end{enumerate}

\begin{enumerate}
\def\labelenumi{\alph{enumi})}
\setcounter{enumi}{1}
\tightlist
\item
  林德勒夫空间与紧空间的积是林德勒夫空间（按第 I 章，§ 9，第 10 号，命题 17 的证明方式论证；参看 § 5，习题 16）。
\end{enumerate}

\begin{enumerate}
\def\labelenumi{\arabic{enumi})}
\setcounter{enumi}{23}
\tightlist
\item
  \begin{enumerate}
  \def\labelenumii{\alph{enumii})}
  \tightlist
  \item
    设 X 是可度量化空间，R 是 X 上的一个闭等价关系，其所有等价类都是紧的。证明 X/R 是可度量化的。{[}对 X 的由关于 R 饱和的集组成的开覆盖证明第 5 号引理 3 的类似命题，以此来应用长田–斯米尔诺夫定理（习题 22）；用 \(F_{n\alpha}\) 表示含于由所有 \(x \in X\) （满足 \(d(x, X - U_{\alpha}) > 2^{-n}\) ）组成之集的最大饱和开集；用 \(F_{n\alpha}'\) 表示由所有 \(x \in X\) （满足 \(d(x, X - U_{\alpha}) \geqslant 2^{-n}\) ）组成之集的饱和化；并用 \(G_{n\alpha}\) 表示由所有 \(x \in F_{n\alpha}\) （满足 \(x \notin F_{n+1,\beta}'\) ，对每个 \(\beta < \alpha\) ）组成的集。{]}
  \end{enumerate}
\end{enumerate}

\begin{enumerate}
\def\labelenumi{\alph{enumi})}
\setcounter{enumi}{1}
\item
  把 a) 的结果推广到如下情形：关系 R 是闭的，且 X/R 的每个点都有可数的基本邻域系{[}利用 § 2 的习题 20 b){]}。
\item
  设 X 是拓扑空间，\((\mathbf{F}_{\alpha})\) 是 X 的一个局部有限闭覆盖。证明：若诸子空间 \(F_{\alpha}\) 的每一个都可度量化，则 X 可度量化（把 X 看作诸 \(F_{\alpha}\) 的拓扑和的一个商空间）。
\item
  仿紧空间（其中每个点都有一个可度量化的邻域）是可度量化的{[}应用 c)；参看第 I 章，§ 9，习题 12 b)，以及第 IX 章，§ 4，习题 25 d)、§ 2，习题 15 与 § 5，习题 15{]}。
\end{enumerate}

\begin{enumerate}
\def\labelenumi{\arabic{enumi})}
\setcounter{enumi}{24}
\tightlist
\item
  豪斯多夫空间 X 称为亚紧的，如果对 X 的任何开覆盖 R，存在 X 的一个加细 R 的逐点有限开覆盖。
\end{enumerate}

\begin{enumerate}
\def\labelenumi{\alph{enumi})}
\item
  证明亚紧空间的每个闭子空间都是亚紧的，并且若亚紧空间 X 的每个开子空间都亚紧，则 X 的每个子空间都亚紧。
\item
  设 X 是豪斯多夫空间，R 是 X 上的一个闭等价关系，其所有等价类都是紧的。证明：若 X/R 亚紧，则 X 亚紧{[}按习题 21 b) 的方式论证{]}。
\item
  证明既亚紧又可数紧（§ 2，习题 14）的拓扑空间是紧的。{[}利用佐恩引理证明每个逐点有限开覆盖都包含一个极小开覆盖，并利用 § 2，习题 14 a)。{]}
\end{enumerate}

\(d)\) 第 I 章，§ 9，习题 12 中定义的局部紧完全正规空间 \(\mathbf{X} = [a, b[\) 不是亚紧的（*）。

\begin{enumerate}
\def\labelenumi{\alph{enumi})}
\setcounter{enumi}{4}
\tightlist
\item
  设 Y 是可度量化空间，A 是 Y 的一个子集，使得 A 与其余集都在 Y 中稠密。用 X 表示给 Y 赋以由 Y 的开集与集 A 生成的拓扑所得的非正则空间。证明 X 是亚紧的。（对每个 \(x \in X\) ，考虑 x 的一个邻域 \(V_{x}\) ，它含于 X 的给定开覆盖 R 的某个集中，使得 \(V_{x} \subset A\) （若 \(x \in A\) ），且使得 \(V_{x}\) 是 x 在 Y 中的一个邻域（若 \(x \notin A\) ）。注意 X 的子空间 A 是仿紧的，而且诸 \(V_{x}\) （\(x \notin A\) ）的并也是仿紧的。）
\end{enumerate}

\begin{enumerate}
\def\labelenumi{\arabic{enumi})}
\setcounter{enumi}{25}
\tightlist
\item
  \begin{enumerate}
  \def\labelenumii{\alph{enumii})}
  \tightlist
  \item
    证明每个伪紧的正规空间 X（§ 1，习题 22）都是可数紧的（§ 2，习题 14）。{[}证明：若 \((x_n)\) 是 X 的互异点组成的、没有触点的一个序列，则 X 上存在一个连续的有限实值函数 f，使得 \(f(x_n) = n\) 对所有 n 成立。{]}
  \end{enumerate}
\end{enumerate}

\begin{enumerate}
\def\labelenumi{\alph{enumi})}
\setcounter{enumi}{1}
\tightlist
\item
  设 \(Y = [a, b[\) 是第 I 章，§ 9，习题 12 中定义的局部紧空间，并设 \(Y_{0}\) 是由 Y 添加一个无穷远点（它可与 b 等同）而得的紧空间。设 c 是 Y 中使 {[}a, c{[} 为无限集的最小元素，并设 Z 是 Y 的紧子空间 {[}a, c{]} 。设 X 是紧积空间 \(Y_{0} \times Z\) 中除去点 (b, c) 后的余集。证明 X 局部紧但非正规{[}参看习题 12 a){]}；X 伪紧但非可数紧。给出 X 的一个非伪紧的闭子空间的例子。再给出 X 上的一个不达到其下确界的下半连续函数的例子{[}参看 § 2，习题 14 i){]}。
\end{enumerate}

\begin{enumerate}
\def\labelenumi{\arabic{enumi})}
\setcounter{enumi}{26}
\tightlist
\item
  拓扑空间 X 称为局部仿紧的，如果 X 的每个点都有一个闭邻域，它是 X 的一个仿紧子空间。
\end{enumerate}

\begin{enumerate}
\def\labelenumi{\alph{enumi})}
\item
  证明局部仿紧空间是完全正则的。
\item
  在习题 26 b) 中定义的紧空间 \(Y_{0} \times Z\) 中，设 H 是由点 \((b, y)\) （其中 y \textless{} c）组成的集的余集。证明 H 正规但非局部仿紧{[}参看第 I 章，§ 9，习题 12 b){]}。
\item
  设 X 是拓扑空间，\(\Phi\) 是 X 的一个由 X 的仿紧子空间闭集组成的覆盖，使得每个 \(A \in \Phi\) 都在 X 中有一个属于 \(\Phi\) 的邻域。设 \(X'\) 是 X 与一个点 \(\omega\) 组成的和集；在 \(X'\) 上如下定义拓扑：取 x 在 X 中的全体邻域所成之集，作为 \(x \in X\) 在 \(X'\) 中的基本邻域系，而 \(\omega\) 在 \(X'\) 中的基本邻域系则取为由 \(\Phi\) 的集的余集所生成的滤子基{[}参看习题 20 c){]}。证明 \(X'\) 是仿紧的。
\end{enumerate}

\begin{enumerate}
\def\labelenumi{\arabic{enumi})}
\setcounter{enumi}{27}
\tightlist
\item
  设 X 是度量空间，d 是 X 上的度量，A 是 X 的一个闭子集，并设 \(f: A \to [1, 2]\) 是连续实值函数。证明 X 上的实值函数 g——\(g(x) = f(x)\) （若 \(x \in A\) ），而
\end{enumerate}

\[
g (x) = \frac {\mathrm{1}}{d (x , \mathrm{A})} \inf _ {y \in \mathrm{A}} (f (y) d (x, y))
\]

若 \(x \in \mathbb{C}A\) 则取上式——在 \(X\) 上连续。由此对度量空间给出定理 2 的一个证明。

\begin{enumerate}
\def\labelenumi{\arabic{enumi})}
\setcounter{enumi}{28}
\item
  设 X 是豪斯多夫拓扑空间。则 X 正规，当且仅当：对 X 的每个闭子集 A 以及每个在 A 上定义且关于诱导拓扑连续的实值函数 f，存在 X 上的一个连续伪度量 d，f 关于它一致连续。（为证条件是充分的，证明每个定义在闭子集 A 上且关于诱导拓扑连续的实值函数 f 都能扩张为 X 上的连续函数；为此，考虑与由伪度量 d 在 X 上定义的一致结构相伴的豪斯多夫空间 \(\overline{X}\) 。）
\item
  设 X 是正规空间，g（相应地，f）是 X 上的上（相应地，下）半连续函数，满足 \(g \leqslant f\) 。证明 X 上存在连续实值函数 h，使得 \(g \leqslant h \leqslant f\) 。{[}化归到 f 与 g 取值于 {[}0, 1{]} 的情形。对每个二进数 \(r \in [0, 1]\) ，设 \(\mathrm{F}(r)\) {[}相应地，\(\mathrm{G}(r)\) {]} 为由所有 \(x \in X\) （满足 \(f(x) \leqslant r\) ，相应地，\(g(x) < r\) ）组成的集。用归纳法定义开集的一个族 \(\mathrm{U}(r)\) （r 跑遍含于 {[}0, 1{]} 的二进数所成的集），使得当 \(r < r'\) 时有 \(\mathrm{F}(r) \subset \mathrm{U}(r')\) 、\(\overline{\mathrm{U}(r)} \subset \mathrm{U}(r')\) 与 \(\overline{\mathrm{U}(r)} \subset \mathrm{G}(r')\) ，办法是模仿第 1 号的定理 1 的证明；然后像定理 1 那样完成证明。{]}
\end{enumerate}
% semantic-fragments: legacy-input-seam
\relax{}
